\documentclass[11pt]{article}
\usepackage[T1]{fontenc}
\usepackage{lmodern,amsmath,amssymb,microtype}
\usepackage[margin=1.02in]{geometry}
\usepackage[colorlinks=true,urlcolor=blue,citecolor=blue]{hyperref}
\setlength{\emergencystretch}{2em}
\hypersetup{pdfauthor={Wentao Li},pdftitle={Dense presentations of the basic Fraenkel model}}
\title{Dense presentations of the basic Fraenkel model}
\author{Wentao Li\\[2pt]\small Prepared with assistance from OpenAI Codex}
\date{November 2026}
\begin{document}
\maketitle
\begin{abstract}
The finite-support permutation model of Fraenkel has a presentation by a group containing no nonidentity internal permutations. Since Frucht's theorem fails there, this refutes Pinsky's proposed converse to his internal-automorphism obstruction. The converse holds for pointwise-closed permutation groups with the finite-support filter.
\end{abstract}

\paragraph{Question.} Start in a choice ground of set theory with atoms. A permutation group $G$ of the atom set $A$ acts on sets by recursion. For the finite-support filter let $M(G)$ consist of the hereditarily symmetric sets: a set has a finite support $E\subseteq A$ if all permutations in the pointwise stabilizer $G_{(E)}$ fix it. In Pinsky's terminology an internal permutation is a member of $G$ whose graph belongs to the model. His Conjecture~13 (arXiv version) says that whenever Frucht's theorem fails in a permutation model, every subgroup in the defining filter contains a nontrivial internal permutation. The filter is allowed to be more general in his question.

\paragraph{Density lemma.} Suppose $D$ is dense in $G$ for agreement on finite subsets of $A$. Then $M(D)=M(G)$ for the respective finite-support filters. Inclusion from right to left follows by restricting the action. For the reverse use induction on rank over atoms. If $x\in M(D)$ has $D$-support $E$, every $y\in x$ is already in $M(G)$ and has some finite $G$-support $F$. Given $h\in G_{(E)}$, choose $g\in D$ agreeing with $h$ on $E\cup F$. Then $hy=gy\in x$. Applying this also to $h^{-1}$ gives $hx=x$. The induction starts with atoms, each supported by its singleton.

\paragraph{A dense group.} Take $A$ countably infinite in the ground. There is a free subgroup $D$ dense in $\operatorname{Sym}(A)$. For completeness, build countably many permutations $g_i$ by finite partial injections. Schedule countably many requirements, including putting each atom in the domain and range of each generator, which can always be done using fresh values. For each finite partial bijection of atoms require that an as-yet untouched generator extend it. For every nonempty reduced word in the $g_i^{\pm1}$ require a moved point. To meet this last requirement take a path on new, pairwise distinct atoms, prescribing the successive generator maps along its edges, inverting the prescription for inverse letters. A collision of sources or targets for one generator along this path would require consecutive inverse letters. Thus the prescriptions are consistent and the endpoints differ. Only finitely many generators and atoms have been touched at each finite stage. The unions are permutations, the finite-bijection requirements give density, and the word requirements give freeness. In particular $D$ is torsion-free.

\paragraph{Counterexample.} In $M(\operatorname{Sym}(A))$ every internal bijection of atoms has finite moved set. Indeed its graph is supported by a finite $E$, so the bijection $f$ commutes with every permutation fixing $E$ pointwise. The fixed points of this stabilizer are precisely $E$, hence $f$ preserves $E$. For $a\notin E$, the stabilizer of $E\cup\{a\}$ fixes $f(a)$, forcing $f(a)=a$. Every permutation with finite moved set has finite order. The density lemma therefore gives $D\cap M(D)=\{1\}$, identifying permutations with their graphs. Yet $M(D)$ is the basic Fraenkel model. In its full symmetric presentation, every filter subgroup contains a nonidentity finitely supported transposition, so Frucht's theorem fails there by Pinsky's Corollary~1. Truth in the model is independent of its presentation. Taking the filter member $D$ itself refutes the conjecture.\hfill$\square$

\paragraph{Where the converse survives.} Suppose $G$ is closed in $\operatorname{Sym}(A)$ for pointwise convergence. If some filter subgroup has no nonidentity internal element, the same is true of $H=G_{(E)}$ for some finite $E$. The group $H$ is closed. In the choice ground, list by an ordinal all its orbits on finite tuples of atoms, recording the arity with each orbit. Each orbit and the ordinal-indexed list are supported by $E$: $H$ preserves each orbit and fixes the pure indices. All entries are hereditarily symmetric, including all individual finite tuples. The relational structure on the atoms with one predicate for each listed orbit is therefore available in the model. Its ground automorphism group is $H$: preserving these predicates gives agreement with an element of $H$ on each finite tuple, which puts a permutation in the closure of $H$. Its internal automorphism group is trivial.

Here are details of converting this internally rigid structure into the rigid graph labels required for Frucht's theorem. Use an ordinal sort with its order containing indices for the predicates and finite positions. Represent true tuples by incidence nodes linked to the predicate index, and their argument occurrences to the tuple, position and atom, using separate binary relation names. The arity and index data are fixed ordinal data. All auxiliary nodes are determined by their incidence, so automorphisms of this finite-signature structure restrict bijectively to the original automorphisms. For each atom $a$, add a unary predicate marking $a$. An internal isomorphism between the structures for $a$ and $b$ restricts to an internal automorphism taking $a$ to $b$, forcing $a=b$.

For clarity a rigid simple graph encoding needs to prevent automorphisms of its own marker gadgets. Make an element node for each element, a node per true relation tuple, and one incidence node per argument position joined to its tuple and its element. Use distinct colours for the roles, relation names and positions. Add a uniquely coloured universal anchor and subdivide each skeleton edge once, giving subdivisions their own colour. At each skeleton vertex attach two odd cycles sharing just this vertex, one of length three and one of length $5+2i$ for colour $i$. In each cycle append a leaf to one specified neighbour of the attachment. The initial subdivided skeleton is bipartite. Its vertices are recovered as the intersections of pairs of odd cycles, and their colours by the longer cycle lengths. Each cycle with its attachment and marked neighbour fixed has only the identity automorphism. Thus restriction and extension give exactly the automorphisms of the original structure. All steps use tagged tuples, finite indices and the previously supplied listing. Point the graph at its anchor. The graphs constructed for atoms are internally rigid and pairwise nonisomorphic, as a single internally defined indexed family. Pinsky's atom-labelling theorem now gives Frucht's theorem in $M(G)$. Together with his obstruction this proves the conjectured criterion for closed finite-support presentations.

\paragraph{Source and priority.} B. Pinsky, \emph{Frucht's theorem without choice}, §5, Conjecture~13, \url{https://arxiv.org/abs/2305.11382v1}. The failure of Frucht's theorem in the basic Fraenkel model and the sufficient condition using graph labels are his results. Our counterexample changes the presenting group and identifies why closedness restores the converse. Independent proof and priority review are pending.
\end{document}
